Write $Q(x)=x\cdot x$ on $\R^3$. A Euclidean lattice is called integral when
all its inner products are integers; its covolume is the volume of a
fundamental parallelepiped. A vector of a lattice is primitive when it is
not in $qL$ for any rational prime $q$.

\begin{lemma}[Small unimodular lattices]\label[lemma]{lem:small-lattices}
Every integral positive definite Euclidean lattice of rank $1\le r\le3$ and
covolume one has an orthonormal $\Z$-basis.
\end{lemma}
\begin{proof}
For $r=3$, the open ball of radius $4/3$ has volume
$256\pi/81>8$. Minkowski's theorem~\cite{minkowski} gives a nonzero lattice vector of
squared length less than $16/9<2$. Its squared length is a positive
integer, so it equals one. The analogous ball in rank two has area
$16\pi/9>4$, and in rank one has length $8/3>2$, so the same reasoning
applies in those ranks.

If $u\in L$ has squared length one, integrality implies
$x\cdot u\in\Z$ for every $x\in L$. Therefore
\[
 L=\Z u\mathbin{\perp}L',\qquad L'=L\cap u^\perp,
\]
because $x-(x\cdot u)u\in L'$. Orthogonal splitting preserves covolume
when the first factor has length one; thus $L'$ is integral of covolume
one. Induction on rank completes the proof.
\end{proof}

\begin{lemma}[Isotropic lines over an odd prime field]\label[lemma]{lem:isotropic-lines}
Let $p$ be an odd prime and let $\mathcal C$ be the set of isotropic
one-dimensional subspaces of $(\F_p^3,Q)$. Then $|\mathcal C|=p+1$.
For nonzero $\bar x\in\F_p^3$, the number of $\ell\in\mathcal C$ with
$\ell\subseteq\bar x^\perp$ is
\[
 1+\left(\frac{-Q(\bar x)}p\right).
\]
For $\bar x=0$, that number is $p+1$.
\end{lemma}
\begin{proof}
There is a nonzero isotropic vector: the sets of squares and of numbers
$-1-y^2$ each have $(p+1)/2$ elements and therefore intersect. This gives
a vector $(x,y,1)$ of norm zero.

Fix a nonzero isotropic vector $u$. Nondegeneracy gives $v$ with
$u\cdot v=1$, and $w=v-Q(v)u/2$ then satisfies
$Q(w)=0$ and $u\cdot w=1$. The span of $u,w$ is nondegenerate;
its one-dimensional orthogonal complement has a generator $e$ with
$Q(e)=c\ne0$. Thus the form in this basis is $2ab+ct^2$.
Replacing $a$ by $a'=-2a/c$ turns its zero equation into $a'b=t^2$.
This conic is parametrized bijectively by
$[u:v]\mapsto[u^2:v^2:uv]$. Explicitly, when its first coordinate
is nonzero, rescale to $(1,s^2,s)$; its only point with first
coordinate zero is $[0:1:0]$. It therefore has $p+1$ points.

If $Q(\bar x)=0$, use the same hyperbolic basis with $u=\bar x$.
Its orthogonal plane is spanned by $u,e$, with norm $ct^2$ on
$au+te$. Thus its sole isotropic line is $\F_p\bar x$.

If $Q(\bar x)\ne0$, the orthogonal plane is nondegenerate. Its Gram
determinant has square class $Q(\bar x)^{-1}$: adjoining $\bar x$ to
any basis of the plane gives Gram determinant $Q(\bar x)$ times the
plane determinant, whereas change of basis in the standard ternary
space changes determinant by a square. A nondegenerate binary form
has two isotropic lines when the negative of its determinant is a
square, and no isotropic lines otherwise. To see this explicitly,
diagonalize it as $As^2+Bt^2$, with $A,B\ne0$ (characteristic is
odd). An isotropic line has $t\ne0$, and $s/t$ must be a root of
$X^2=-B/A$. There are two or zero roots according as the negative
determinant $-AB$ is a square or a nonsquare. The square classes of
$-Q(\bar x)^{-1}$ and $-Q(\bar x)$ coincide. The asserted expression
follows in both cases.
\end{proof}

\begin{lemma}[The neighbors of $\Z^3$]\label[lemma]{lem:neighbors}
Put $L=\Z^3$. For each $\ell\in\mathcal C$ there is an integral
covolume-one lattice $L_\ell$ satisfying
\[
 \begin{gathered}
 L\cap L_\ell=M_\ell
    :=\{x\in L:x\bmod p\in\ell^\perp\},\\
 [L:M_\ell]=[L_\ell:M_\ell]=p,\\
 pL\subseteq M_\ell,\qquad pL_\ell\subseteq M_\ell\subseteq L.
 \end{gathered}
\]
If $v\in L\setminus pL$ and $p^2\mid Q(v)$, then $v/p$ belongs
to exactly one of these lattices, namely $L_{\F_p(v\bmod p)}$.
\end{lemma}
\begin{proof}
Choose a lift $u\in L$ of a nonzero vector of $\ell$. It can be chosen
with $p^2\mid Q(u)$. Indeed, starting with any lift, replacing it by
$u+pz$ changes $Q(u)/p$ modulo $p$ by $2u\cdot z$. The linear
functional $z\mapsto u\cdot z\pmod p$ is nonzero, hence surjective,
so a correction $z$ with the required value exists.

Set
\[
 L_\ell=M_\ell+\Z\frac up.
\]
All inner products are integral, since $u\cdot M_\ell\subseteq p\Z$
and $Q(u)\in p^2\Z$. The functional defining $M_\ell$ is surjective
modulo $p$, so its index in $L$ is $p$. The class of $u/p$ modulo
$M_\ell$ has order $p$, since $u\notin pL$ and $u\in M_\ell$.
Consequently $L_\ell$ has covolume one. If $m+tu/p\in L$, then
$p\mid t$, proving the intersection assertion. All other inclusions
follow from the definition. Each $L_\ell$ is isometric as a lattice
to $\Z^3$ by the preceding small-lattice lemma.

For the last assertion, membership forces the line of $v\bmod p$
to equal $\ell$. Conversely write $v=tu+pz$ with $p\nmid t$.
The divisibility of $Q(v)$ and $Q(u)$ by $p^2$ gives
$2t(u\cdot z)\equiv0\pmod p$. Thus $z\in M_\ell$ and
$v/p=tu/p+z\in L_\ell$.
\end{proof}

The construction is independent of the normalized lift: if another lift
is $tu+pz$ with $t\not\equiv0\pmod p$ and norm divisible by $p^2$,
the calculation in the last paragraph shows $z\in M_\ell$; the two
generators then generate the same cyclic extension of $M_\ell$.

\paragraph{Primitivity away from $p$.}
In the double count below, passage between $L$ and $L_\ell$ preserves
primitivity at every prime $q\ne p$. This can be checked without using
localization. A common vector $x\in M_\ell$ is in $qL$ if and only
if it is in $qL_\ell$. Indeed, if $x=qz$ with $z\in L$, reduction
modulo $p$ shows $z\in M_\ell$. If instead $z\in L_\ell$, its
class in the group $L_\ell/M_\ell$ of order $p$ is annihilated by
$q$ and hence is zero; again $z\in M_\ell$. For an
outside vector $y=v/p\in L_\ell\setminus L$, the corresponding
statement is $y\in qL_\ell$ if and only if $v\in qL$. In the latter
direction, write $v=qz$. Since $v\in M_\ell$ and $q$ is invertible
modulo $p$, $z\in M_\ell$. Choose integers $a,b$ with $aq+bp=1$;
then $z/p=ay+bz\in L_\ell$. The converse follows by multiplying by
$p$, since $pL_\ell\subseteq L$.

Recall that $F(m)$ is the number of primitive vectors of norm $3m^2$ in
$L$. These sets are finite by coordinate bounds. All the lattices
$L_\ell$ have exactly $F(m)$ such vectors, because they are
isometric to $L$ by a bijective $\Z$-linear isometry. Let
$\chi_p=\bigl(\frac{-3}{p}\bigr)$.

\begin{lemma}[Primitive double count]\label[lemma]{lem:double-count}
Let $p$ be an odd prime and let $m$ be positive and odd.
If $p\nmid m$, then
\[
 F(pm)=(p-\chi_p)F(m).
\]
If $p\mid m$, then
\[
 F(pm)=(p+1)F(m)-
 \left(p-\left(\frac{-3(m/p)^2}{p}\right)\right)F(m/p).
\]
\end{lemma}
\begin{proof}
Count pairs $(\ell,y)$ where $y$ is primitive in $L_\ell$ and has
norm $3m^2$. Their total number is $(p+1)F(m)$.

First count the pairs with $y\notin L$. The vector $v=py$ is in
$L\setminus pL$ and has norm $3p^2m^2$. Conversely every primitive
such $v$ gives exactly one outside pair by the neighbor lemma.
Primitivity away from $p$ is preserved as just proved. An outside
$y$ is automatically primitive at $p$, because $y\in pL_\ell$
would imply $v\in p^2L_\ell\subseteq pL$. Thus the outside pairs
number exactly $F(pm)$.

Suppose first that $p\nmid m$. The integer $3m^2$ is not divisible
by $p^2$, including at $p=3$. Therefore every vector of this norm
in an integral lattice is primitive at $p$. Inside pairs are precisely
primitive $x\in L$ with $\ell\subseteq(x\bmod p)^\perp$.
The finite-field lemma gives $1+\chi_p$ choices of $\ell$ for each
such $x$: for $p\ne3$ the factor $m^2$ is a nonzero square, and for
$p=3$ both Legendre symbols are zero. Therefore
$(p+1)F(m)=F(pm)+(1+\chi_p)F(m)$, giving the first formula.

Now suppose $p\mid m$. An inside vector that is primitive away
from $p$ has the unique form $x=p^az$ with $a\ge0$ and $z$
primitive in $L$. If $a=0$, then $x\bmod p$ is nonzero and
isotropic. Its unique orthogonal isotropic line is its own line;
for this neighbor $x/p\in L_\ell$ by the last part of the neighbor
lemma. Hence $x$ is not primitive in that neighbor and contributes
nothing. If $a\ge2$, then $x/p\in pL\subseteq L_\ell$ for every
neighbor, so again it contributes nothing.

Only $a=1$ contributes. Write $x=pz$, with $z$ primitive of norm
$3(m/p)^2$. The vector $x$ belongs to all $p+1$ neighbors. It fails
to be primitive at $p$ precisely when $z\in L_\ell$, or equivalently
when $\ell\subseteq(z\bmod p)^\perp$. The finite-field lemma shows
that the number of remaining choices is
$p-\bigl(\frac{-3(m/p)^2}{p}\bigr)$. Counting these pairs yields
the second formula.
\end{proof}

\begin{theorem}[Primitive three-square product]\label[theorem]{thm:three-squares}
For every positive odd $m$,
\[
 F(m)=8m\prod_{p\mid m}\left(1-\frac{\chi_p}{p}\right).
\]
\end{theorem}
\begin{proof}
First $F(1)=8$, since the vectors of norm three are precisely
$(\pm1,\pm1,\pm1)$.

Fix an odd prime $p$ and a positive odd $b$ with $p\nmid b$.
The first double-counting identity gives
$F(pb)=(p-\chi_p)F(b)$. Apply the second identity at $m=pb$.
Its Legendre symbol is $\chi_p$, so
\[
 F(p^2b)=(p+1)(p-\chi_p)F(b)-(p-\chi_p)F(b)
        =p(p-\chi_p)F(b).
\]
For exponent $a\ge2$, the same identity at $m=p^ab$ has zero
Legendre symbol, since $p\mid p^{a-1}b$. Induction therefore gives
\[
 F(p^{a+1}b)=(p+1)F(p^ab)-pF(p^{a-1}b)
            =p^a(p-\chi_p)F(b).
\]
Removing the prime-power factors of $m$ one at a time now gives
\[
 F(m)=8\prod_{p^a\parallel m}p^{a-1}(p-\chi_p)
     =8m\prod_{p\mid m}(1-\chi_p/p).
\]
The local symbol agrees with the mod-six character by quadratic
reciprocity, as proved explicitly in Section~\ref{sec:character}.
\end{proof}
\paragraph{The prime three.}
No use has been made of $p\nmid3$ in the neighbor argument. The
underlying Gram matrix of $L$ is the identity and stays nondegenerate
modulo every odd prime, including three. If $3\nmid m$, the norm
$3m^2$ has valuation exactly one, so inside vectors are still
primitive at three and the isotropic-line count is one. This gives
$F(3m)=3F(m)$. If $3\mid m$, the symbol in the second recurrence is
always zero; it gives
$F(3m)=4F(m)-3F(m/3)$, whose induction again yields $3F(m)$.
Thus the prime three contributes its full power through $m$, while
its Euler factor is one.

